% !PRINTED 51
\weberpage{51}

Es ist vielleicht erwünscht, diese Sätze an einem einfachen
Beispiele zu erläutern. Es sei der Grad der Gruppe $S$ das Quadrat
einer Primzahl $n = p^2$. Dann hat $S$ entweder nur die eine
Invariante $p^2$ und ist dann cyklisch:
\[
\text{a)}\quad S = 1, A, A^2 \dots A^{p^2-1},
\]
oder es hat $S$ zwei Invarianten, die beide gleich $p$ sind; dann
besteht $S$ aus den Elementen:
\[
\text{b)}\quad S = A_1^{\alpha_1}A_2^{\alpha_2}, \quad \alpha_1,
\alpha_2 = 0, 1 \dots p-1.
\]
Im Falle a) erhalten wir die $p^2$ Charaktere
\[
\chi_B(A^\alpha) = \varepsilon^{\beta\alpha},
\]
wenn $\varepsilon$ eine primitive Wurzel der Gleichung
$\varepsilon^{pp} = 1$ ist.
% [sic: the source prints \varepsilon^{pp}; the cyclic group order and surrounding p^2 formulas require \varepsilon^{p^2}.]
Nehmen wir, um nach dem Satze 10. die
Theiler von $S$ zu bilden, einen der Charaktere $\chi_B$, in dem
$\beta$ nicht durch $p$ theilbar ist, so wird $\chi_B(A^\alpha)$
nur dann $= 1$ sein können, wenn $\alpha$ durch $p^2$ theilbar
oder also $= 0$ ist, d.\,h. wir bekommen nur den aus dem
Einheitselemente bestehenden Theiler von $S$. Nehmen wir aber
$\beta = p\beta'$ durch $p$, aber nicht durch $p^2$ theilbar an, so
wird $\chi_{p\beta'}(A^\alpha)$ dann und nur dann $= 1$, wenn
$\alpha = p\alpha'$ durch $p$ theilbar ist. Wir erhalten also den
Theiler
\[
T = 1, A^p, A^{2p} \dots A^{(p-1)p},
\]
und der zu $T$ reciproke Theiler $U$ ist hier mit $T$ identisch.

Im Falle b) müssen wir, um die Charaktere zu bilden, zwei
$p^{\text{te}}$ Einheitswurzeln $\varepsilon^{\beta_1},
\varepsilon^{\beta_2}$ annehmen, und erhalten, wenn
\[
B = A_1^{\beta_1}A_2^{\beta_2}
\]
gesetzt ist,
\[
\chi_B(A_1^{\alpha_1}A_2^{\alpha_2}) = \varepsilon^{\alpha_1\beta_1 + \alpha_2\beta_2}.
\]
Setzen wir zwei dieser Charaktere $= 1$, also
\[
\alpha_1\beta_1 + \alpha_2\beta_2 \equiv 0, \quad
\alpha_1\beta_1' + \alpha_2\beta_2' \equiv 0,
\]
so folgt, wenn die Determinante $\beta_1\beta_2' - \beta_2\beta_1'$
nicht durch $p$ theilbar ist, dass $\alpha_1$ und $\alpha_2$ durch
$p$ theilbar sein müssen. Ist aber die Determinante durch $p$
theilbar, so ist die eine dieser beiden Congruenzen eine Folge der
anderen. Im ersten Falle bekommen wir also eine Gruppe, die nur
aus dem Einheitselemente besteht. Wir erhalten daher alle von 1
und $S$ verschiedenen Theiler $T_{\beta_1,\beta_2}$, wenn wir für
ein feststehendes $\beta_1, \beta_2$ die $\alpha_1, \alpha_2$ auf
alle möglichen Arten der Congruenz
\weq{9}{\alpha_1\beta_1 + \alpha_2\beta_2 \equiv 0\ (\mathrm{mod}\ p)}
%CONT>>
